%%%%%%%%%%%%%%%%%%%%%%%%
%
% $Autor: Wings $
% $Datum: 2019-07-09 09:26:07Z $
% $Pfad: Vorlesungen/WS_19_20/Projekte/KandaNeuralnetwork/latex - Ausarbeitung/Kapitel/Hardware.tex $
% $Version: 4440 $
%
%%%%%%%%%%%%%%%%%%%%%%%%



\chapter{Google Colab}

Google Colaboratory oder Colab ist ein Produkt von Google Research und ist eine kostenlose Jupyter-Notebook-Umgebung, die vollständig in der Cloud läuft. Es ist in der Regel gut für maschinelles Lernen und Datenanalyse geeignet. Das Gute daran ist, dass es keine Einrichtung erfordert und freien Zugang zu Rechenressourcen wie GPU und TPU bietet. \cite{GoogleColab:2021}.

\bigskip

\textbf{Vorteile von Colab} 

\begin{itemize}  
	\item \textbf{Vorinstallierte Bibliotheken:}
	
	Es bietet vorinstallierte Bibliotheken für maschinelles Lernen, einschließlich PyTorch, Keras, TensorFlow
	\item \textbf{Speichern in der Cloud:}
	
	Wenn wir ein Jupyter-Notebook verwenden, wird alles auf einem lokalen Gerät gespeichert. Wenn wir jedoch Google Colab verwenden, können wir es in der Cloud speichern und von jedem Gerät aus darauf zugreifen, indem wir uns einfach bei unserem Google Drive-Konto anmelden.
	\item \textbf{Collaboration:}
	
	Es hilft dabei, mehreren Entwicklern, die an einem Projekt arbeiten, Zugang zu gewähren, um den Code gemeinsam zu bearbeiten und den fertigen Code mit den Entwicklern zu teilen.
	\item \textbf{Freie GPU- und TPU-Nutzung:}  
	
	Dies ist wahrscheinlich die beste Funktion, da Google Research uns die Nutzung seiner dedizierten GPUs und TPUs für unsere persönlichen Machine-Learning-Projekte ermöglicht. Daher müssen unsere Computer beim Training des Modells keine komplexen Berechnungen durchführen, da dies in der Cloud geschieht.
	
\end{itemize}

Google Colaboratory, oder kurz \glqq Colab\grqq, ist ein Produkt von Google Research. Colab ermöglicht es jedem, beliebigen Python-Code im Browser zu schreiben und auszuführen, und eignet sich besonders gut für maschinelles Lernen, Datenanalyse und Bildung.
Mit Colab können wir Python im Browser schreiben und ausführen, mit

\begin{itemize}
	\item Null Konfiguration erforderlich
	\item Freier Zugang zu GPUs
	\item Einfache Freigabe
\end{itemize}

\section{Notebook und Zellen}

Colab bietet ein interaktives Notizbuch, das eine Liste von Zellen enthält, die entweder erklärenden Text oder ausführbaren Code und dessen Ausgabe enthalten. Klicken Sie auf eine Zelle, um sie auszuwählen.
Die Textzellen werden mit einer einfachen Auszeichnungssprache namens Markdown formatiert.
Um den Markdown-Quelltext zu sehen, doppelklicken Sie auf eine Textzelle, um sowohl den Markdown-Quelltext als auch die gerenderte Version anzuzeigen. Über der Markdown-Quelle befindet sich eine Symbolleiste zur Unterstützung der Bearbeitung.

Klicken Sie in die Zelle, um sie auszuwählen, und führen Sie den Inhalt auf die folgenden Arten aus:
Im Folgenden sind einige der Operationen aufgeführt, die zur Ausführung einer Zelle durchgeführt werden können: 

\begin{itemize}
	\item Klicken Sie auf das Abspielen-Symbol \glqq Play icon \grqq{} in der linken Seitenrinne der Zelle;
	\item \glqq Cmd/Ctrl+Enter\grqq{} , um die Zelle an Ort und Stelle auszuführen;
	\item \glqq Shift+Enter\grqq{} , um die Zelle auszuführen und den Fokus auf die nächste Zelle zu verschieben (und eine hinzuzufügen, wenn keine vorhanden ist); oder
	\item \glqq Alt+Enter\grqq{} , um die Zelle auszuführen und eine neue Codezelle direkt darunter einzufügen.
\end{itemize} 

Es gibt zusätzliche Optionen für die Ausführung einiger oder aller Zellen im Menü \glqq Runtime\grqq{} (siehe unten). 

\begin{figure}[H]
	\centering
	\GRAPHICS{1}{1}{GoogleColab/cell}
	\caption{Ein Beispiel für eine Zelle in einem Google Colab Notebook}
\end{figure}

Es gibt auch Textzellen, in denen wir Titel oder Details des Codes hinzufügen können. 
Sie können neue Zellen hinzufügen, indem Sie die Schaltflächen + CODE und + TEXT verwenden, die angezeigt werden, wenn Sie den Mauszeiger über die Zellen bewegen. Diese Schaltflächen befinden sich auch in der Symbolleiste oberhalb des Notizbuchs, wo sie verwendet werden können, um eine Zelle unterhalb der aktuell ausgewählten Zelle hinzuzufügen.
Sie können eine Zelle verschieben, indem Sie sie auswählen und in der oberen Symbolleiste auf Zelle nach oben oder Zelle nach unten klicken.
Aufeinanderfolgende Zellen können durch "Lasso-Auswahl" ausgewählt werden, indem Sie von außerhalb einer Zelle durch die Gruppe ziehen. Nicht benachbarte Zellen können gleichzeitig ausgewählt werden, indem Sie auf eine Zelle klicken und dann die Strg-Taste gedrückt halten, während Sie auf eine andere klicken. Wenn Sie die Umschalttaste anstelle der Strg-Taste drücken, werden alle dazwischen liegenden Zellen ausgewählt.

\section{Google Colab und Machine Learning}

Mit Colab können Sie einen Bilddatensatz importieren, einen Bildklassifikator darauf trainieren und das Modell auswerten - und das alles in nur wenigen Codezeilen. Colab-Notebooks führen den Code auf den Cloud-Servern von Google aus. Das bedeutet, dass Sie die Leistung der Google-Hardware, einschließlich GPUs und TPUs, unabhängig von der Leistung Ihres Computers nutzen können. Alles, was Sie brauchen, ist ein Browser.
Colab wird in der Community für maschinelles Lernen ausgiebig genutzt, unter anderem für folgende Anwendungen:

\begin{itemize}
	\item Erste Schritte mit TensorFlow
	\item Entwickeln und Trainieren neuronaler Netze
	\item Experimentieren mit TPUs
	\item KI-Forschung verbreiten
	\item Tutorials erstellen
\end{itemize}	

\section{Integration mit Drive}

Colaboratory ist mit Google Drive integriert. Es ermöglicht Ihnen, dasselbe Dokument mit mehreren Personen zu teilen, zu kommentieren und gemeinsam daran zu arbeiten:

\begin{itemize}
	\item Über die Schaltfläche TEILEN (oben rechts in der Symbolleiste) können Sie das Notizbuch freigeben und die dafür festgelegten Berechtigungen kontrollieren.
	\item Datei->Kopie erstellen erstellt eine Kopie des Notizbuchs in Drive.
	\item Der Menüpunkt Datei->Speichern speichert die Datei im Laufwerk. Datei->Speichern und Kontrollpunkt fixiert die Version, damit sie nicht aus dem Revisionsverlauf gelöscht wird.
	\item Der Punkt Datei->Revisionsverlauf zeigt den Revisionsverlauf des Notizbuchs an.
\end{itemize}

\section{Verwendung von Colab mit Github}

Das Laden eines Notizbuchs aus einem privaten GitHub-Repository ist möglich, erfordert aber einen zusätzlichen Schritt, damit Colab auf Ihre Dateien zugreifen kann. Gehen Sie wie folgt vor:

\begin{itemize}
	\item Navigate to \SHELL{http://colab.research.google.com/github}.
	\item Klicken Sie auf das Kontrollkästchen \glqq Include Private Repos\grqq.
	\item Melden Sie sich in dem Popup-Fenster bei Ihrem Github-Konto an und autorisieren Sie Colab, die privaten Dateien zu lesen.
	\item Ihre privaten Repositories und Notebooks sind nun über den Github-Navigationsbereich verfügbar.
\end{itemize}

\subsection{Speichern von Notizbüchern auf GitHub oder Drive}

Jedes Mal, wenn Sie ein auf GitHub gehostetes Notizbuch in Colab öffnen, wird eine neue bearbeitbare Ansicht des Notizbuchs geöffnet. Sie können das Notizbuch ausführen und ändern, ohne sich Gedanken über das Überschreiben des Quellcodes zu machen.
Wenn Sie Ihre Änderungen in Colab speichern möchten, können Sie das geänderte Notizbuch über das Menü "Datei" entweder auf Google Drive oder zurück auf GitHub speichern. Wählen Sie File→Save a copy in Drive oder File→Save a copy to GitHub und folgen Sie den Anweisungen. Um ein Colab-Notizbuch auf GitHub zu speichern, müssen Sie Colab die Erlaubnis erteilen, den Commit in Ihr Repository zu übertragen.

\section{Wie man einen Code in Colab ausführt}

Das Ausführen von Code in Google Colab ist so einfach wie das Öffnen einer beliebigen Website. Es sind nur 2 Schritte erforderlich: Melden Sie sich bei Google Colab an und erstellen Sie ein neues Notizbuch.
Sie können ein neues Notizbuch erstellen, indem Sie auf das DAtei-Menü  gehen und \glqq New notebook\grqq{} auswählen (siehe Abbildung unten), 

\begin{figure}[H]
	\centering
	\GRAPHICS{1}{1}{GoogleColab/notebook}
	\caption{Ein neues Colab-Notizbuch mit verschiedenen Optionen}
\end{figure}

Sie können den Namen des Notizbuchs ändern, indem Sie oben links neben dem Google Drive-Logo auf den Dateinamen doppelklicken. Ändern Sie jedoch nicht die Erweiterung der Datei. Das Notizbuch sollte immer die Erweiterung \glqq ipynb\grqq{} haben. Sie können nun Ihren Python-Code schreiben. Auf der Hauptseite gibt es zwei Optionen, um Ihren Code oder Text zu schreiben: Text ist für die Information oder die Beschreibung des Codes. Er ist nur zur Anzeige gedacht. Code ist der Ort, an dem du deinen Code schreibst.

Klicken Sie auf die Schaltfläche \glqq Play\grqq{} auf der linken Seite des Code-Editors, um Ihr Programm auszuführen. Unten sehen Sie den Play-Button für eine Zelle. 

\begin{figure}[H]
	\centering
	\GRAPHICS{1}{1}{GoogleColab/play}
	\caption{Abspieltaste für eine Zelle, um sie auszuführen}
\end{figure}

Standardmäßig haben die Notebooks eine maximale Lebensdauer von bis zu 12 Stunden. Es wird daher empfohlen, Ihre Arbeit zu speichern, bevor Sie Colab verlassen. Es gibt viele Möglichkeiten, Ihr Notizbuch zu speichern. Wie beim Öffnen eines neuen Notizbuchs können die Notizbücher auf Ihrem Google Drive, im lokalen Speicher oder in Ihren Git-Repositories gespeichert werden.
Standardmäßig sind die GPU und TPU für die Notizbücher deaktiviert. Sie können jedoch leicht eingeschaltet werden. Die Schritte zum Einschalten der GPU sind:

\begin{itemize}
	\item Klicken Sie auf Runtime in der Menüleiste.
	\item Wählen Sie die Option Laufzeittyp ändern.
	\item Wählen Sie GPU/TPU als Hardwarebeschleuniger aus.	
\end{itemize}

Unten sehen Sie ein Bild der Einstellungen für den Hardware-Beschleuniger: 

\begin{figure}[H]
	\centering
	\GRAPHICS{1}{1}{GoogleColab/runtime}
	\caption{Harware Accelerator-Einstellungen zur Ausführung eines Programms}
\end{figure}


\section{Google Colaboratory und TensorFlow-Framework} 

TensorFlow ist ein maschinelles Lernsystem, das in großem Maßstab und in heterogenen Umgebungen arbeitet. Es verwendet Datenflussgraphen, um Berechnungen, gemeinsame Zustände und die Operationen, die diesen Zustand verändern, darzustellen. TensorFlow kann tiefe neuronale Netze für die Klassifizierung handgeschriebener Ziffern, Bilderkennung, Worteinbettungen, rekurrente neuronale Netze, Sequenz-zu-Sequenz-Modelle für die maschinelle Übersetzung, die Verarbeitung natürlicher Sprache und PDE (partielle Differentialgleichungen) basierte Simulationen trainieren und ausführen. Es bildet die Knoten eines Datenflussgraphen über viele Maschinen in einem Cluster und innerhalb einer Maschine über mehrere Recheneinheiten ab, einschließlich Multicore-CPUs, Allzweck-GPUs und kundenspezifischer ASICs, bekannt als Tensor Processing Units (TPUs).TensorFlow ermöglicht es Entwicklern, mit neuen Optimierungen und Trainingsalgorithmen zu experimentieren. TensorFlow unterstützt eine Vielzahl von Anwendungen, wobei der Schwerpunkt auf dem Training und der Inferenz auf tiefen neuronalen Netzen liegt. Mehrere Google-Dienste verwenden TensorFlow in der Produktion, wir haben es als Open-Source-Projekt veröffentlicht, und es ist in der Forschung für maschinelles Lernen weit verbreitet. In diesem Papier beschreiben wir das TensorFlow-Datenflussmodell und demonstrieren die überzeugende Leistung, die TensorFlow für mehrere reale Anwendungen erreicht.\href{https://www.tensorflow.org/overview/}{TensorFlow}

\subsection{TensorFlow-Anwendungsfall}

TensorFlow ist ein von Google entwickeltes Framework zur Erstellung von Anwendungen für maschinelles Lernen und zum Training des Modells. Die Google-Suchmaschine basiert ebenfalls auf der KI-Anwendung TensorFlow, z. B. wenn ein Benutzer \glqq a\grqq{} in die Google-Suchleiste eingibt, sagt die Suchmaschine das am besten geeignete vollständige Wort voraus, das ausgewählt werden soll, oder sie kann Vorhersagen treffen, die von der vorherigen Suche des Benutzers im selben System abhängen. TensorFlow kann für eine Reihe von Aufgaben verwendet werden, hat aber einen besonderen Fokus auf das Training und die Inferenz von tiefen neuronalen Netzen. TensorFlow verfügt über eine Reihe von Arbeitsabläufen zur Entwicklung und zum Training des Modells mit den Programmiersprachen Python und Javascripts. Nach dem Training des Modells können wir das Modell in der Cloud oder auch auf einem beliebigen Gerät für die Erstellung von Edge-Computing-Anwendungen einsetzen, unabhängig davon, welche Sprache im Gerät verwendet wird. Es liegt auf der Hand, dass der wichtigste Teil jeder Anwendung des maschinellen Lernens darin besteht, das ML-Modell zu trainieren, damit das Modell in Echtzeit angewendet werden kann und die Entscheidung ohne menschliches Eingreifen getroffen werden kann, wie es der Mensch normalerweise tut. 

\subsection{TensorFlow Lite für Mikrocontroller}

TFLite (TensorFlow Lite) ist eine Reihe von Werkzeugen zur Konvertierung und Optimierung von TensorFlow-Modellen für den Einsatz auf mobilen und Edge-Geräten.TensorFlow Lite, als eine Edge-KI-Implementierung, beseitigt die Herausforderungen der Integration von groß angelegter Computer Vision mit maschinellem Lernen auf dem Gerät, so dass maschinelles Lernen überall eingesetzt werden kann.TensorFlow Lite für Mikrocontroller ist ein Framework für maschinelles Lernen, das für den Einsatz auf Mikrocontrollern und anderen Geräten mit begrenztem Speicher entwickelt wurde. Wir können die Intelligenz von Milliarden von Geräten in unserem Leben, einschließlich Haushaltsgeräten und Geräten des Internets der Dinge, verbessern, indem wir maschinelles Lernen auf winzige Mikrocontroller bringen, anstatt auf teure Hardware oder zuverlässige Internetverbindungen zu setzen, die oft Bandbreiten- und Energiebeschränkungen unterliegen und zu hohen Latenzzeiten führen. TensorFlow Lite ist speziell für maschinelles Lernen auf dem Gerät (Edge ML) optimiert. Als Edge ML-Modell ist es für den Einsatz auf ressourcenbeschränkten Edge-Geräten geeignet. \cite{GoogleTensorFlowLite:2021,Dokic:2020}



%%%%%%%%%%%%%%%%%%%%%%%%%%%%%%%
\section{Umgebung für die Ausbildung des Modells: Google Colab}
%%%%%%%%%%%%%%%%%%%%%%%%%%%%%%%

Google Colaboratory, oder kurz \glqq Colab\grq, ist ein Produkt von Google Research. Colab ermöglicht es jedem, beliebigen Python-Code im Browser zu schreiben und auszuführen, und eignet sich besonders gut für maschinelles Lernen, Datenanalyse und Bildung.
Colab ermöglicht es uns, Python im Browser zu schreiben und auszuführen, mit folgenden Eigenschaften

\begin{itemize}
	\item Keine Konfiguration erforderlich
	\item Freier Zugang zu GPUs
	\item Einfaches Teilen
\end{itemize}

Weitere Informationen über das Google Colab, wie z.B. eine detaillierte Beschreibung, wie man ein Programm ausführt, Google Drive und Github integriert, finden Sie in der Datei \FILE{colab.tex} im Software-Ordner.

\section{Wahl des Modells: YOLO V4} 

\ac{yolo} ist ein hochmoderner Objektdetektor, der die Objekterkennung in Echtzeit mit einer guten Genauigkeit durchführen kann und zur Familie der einstufigen Erkennung gehört, die auch als One-Shot-Erkennung bezeichnet wird.

\ac{yolo} V4 basiert ebenfalls auf dem Darknet und hat auf dem Datensatz \ac{coco} einen \ac{ap}-Wert von 43,5 Prozent sowie eine Echtzeitgeschwindigkeit von 65 \ac{fps} auf dem Tesla V100 erzielt und damit die schnellsten und genauesten Detektoren sowohl in Bezug auf die Geschwindigkeit als auch die Genauigkeit übertroffen.

Im Vergleich zum \ac{yolo} V3 haben sich \ac{ap} und \ac{fps} um 10 Prozent bzw. 12 Prozent erhöht.

\textbf{Anmerkung:} Weitere Informationen über das Modell \ac{yolo} V4, wie z. B. die detaillierte Architektur, die verschiedenen Schichten und die Hyperparameter, finden Sie in der Datei \FILE{YOLOV4.tex} im Ordner \FILE{Software}.

\section{Training des Modells}

Im Folgenden werden die Schritte zum Trainieren des Modells mit Google Colab beschrieben.

\begin{itemize}
  \item Erstellen von yolov4- und Schulungsordnern im Google Drive
  \item Laufwerk mounten, den Ordner verlinken und zum Ordner yolov4 navigieren
  \item Klone das Darknet git repository \cite{Bochkovskiy:2020}
  \item Erstellen \& Hochladen der Dateien, die wir für das Training benötigen (d.h. \FILE{obj.zip} , \FILE{yolov4- custom.cfg}, \FILE{obj.data}, \FILE{obj.names} und \FILE{process.py} ) auf das Laufwerk
  \item Änderungen im Makefile vornehmen, um OPENCV und \ac{gpu} zu aktivieren
  \item Führen Sie den Befehl make aus, um darknet zu erstellen
        Kopieren Sie die Dateien \FILE{obj.zip}, \FILE{yolov4-custom.cfG} \FILE{obj.data} , \FILE{obj.names} aus dem yolov4-Ordner in das darknet-Verzeichnis
  \item Führen Sie das process.py Py-Skript aus, um die \FILE{train.txt} zu erstellen \& \FILE{test.txt} zu erstellen
  \item Laden Sie die vortrainierten YOLOv4-Gewichte herunter
  \item Den Detektor trainieren
  \item Leistung prüfen
\end{itemize}

\subsection{Architektur von YOLO V4} 

Dieser Abschnitt gibt einen kurzen Überblick über das \ac{yolo} V4-Modell.\ref{ColabOneShot}
Das \ac{yolo} V4 \cite{Bochkovskiy:2020} besteht aus:

\begin{itemize}
  \item Backbone : CSP Darknet53.

	    Backbone ist ein tiefes neuronales Netz, das hauptsächlich aus Faltungsschichten besteht und dessen Hauptziel darin besteht, die wesentlichen Merkmale zu extrahieren. \ac{yolo} V4 wendet das \ac{csp} auf Darknet 53 an, das 53 Faltungsschichten verwendet.

  \item Neck: \ac{spp} und \ac{pan}

        Der Hals wird immer zwischen dem Rückgrat und dem Kopf hinzugefügt und kann Objekte in verschiedenen Maßstäben und räumlichen Auflösungen erkennen. \ac{spp} wendet maximales Pooling an, während \ac{pan} angewendet wird, um Merkmale in einer hierarchischen Struktur anzuziehen.

  \item Head: \ac{yolo} V3

        Der Kopf eines Objektdetektors ist für die Klassifizierung und Lokalisierung zuständig.
\end{itemize}

\begin{figure}[H]
	\centering
	\GRAPHICS{1.0}{1.0}{GoogleColab/oneshotdetect.png}
	\caption{Architektur eines Einzelschussdetektors \cite{Bochkovskiy:2020}}
	\label{ColabOneShot}
\end{figure}

\subsection{Aktivieren der GPU in Google Colab}

Im Colab Notebook kann der \ac{gpu} aktiviert werden, indem man auf Runtime klickt und den Runtime-Typ, Hardwarebeschleuniger auf GPU ändert. Durch die Verwendung des Grafikprozessors kann das YOLO V4 wesentlich schnellere Vorhersagen treffen als der \ac{cpu}. 

\subsection{Aufbau des Darknet}

Darknet ist ein Open-Source-Framework für neuronale Netze, das in C und \ac{cuda} geschrieben wurde. \cite{Bochkovskiy:2020} Es ist schnell, einfach zu installieren und unterstützt CPU- und GPU-Berechnungen.mDas Repository kann mit dem unten stehenden Befehl von Github geklont werden.

\medskip

\SHELL{!git clone https://github.com/AlexeyAB/darknet}

\medskip

Der nächste Schritt besteht darin, zum Ordner des Darknets zu navigieren und die notwendigen Änderungen vorzunehmen, um die GPU und OpenCV\index{OpenCV} zu aktivieren. \ref{ColabGPU}

\begin{figure}[H]
	\centering
	\GRAPHICS{0.8}{0.8}{GoogleColab/makefile.png}
	\caption{GPU und OpenCV im Makefile auf 1 setzen}\label{ColabGPU}
\end{figure}

Die CUDA Version kann mit folgendem Befehl überprüft werden, siehe Abbildung \ref{ColabCUDA}:

\medskip

\SHELL{!/usr/local/cuda/bin/nvcc --version}

\medskip

\begin{figure}[H]
	\centering
	\GRAPHICS{0.8}{0.8}{GoogleColab/cudaversion.png}
	\caption{Verifizierung der CUDA-Version}\label{ColabCUDA}
\end{figure}

Der nächste Schritt ist die Ausführung von \SHELL{!make}.
Mit diesem Befehl wird darknet erstellt, so dass Sie anschließend die ausführbare darknet-Datei verwenden können, um Objektdetektoren auszuführen oder zu trainieren.
Die erfolgreiche Kompilierung des obigen Befehls kann überprüft werden, indem die Datei \FILE{darknet} ohne Erweiterungen im Ordner \glqq darknet\grqq{} überprüft wird, die durch diesen Befehl erstellt wird. \ref{ColabDarknet}

\begin{figure}[H]
	\centering
	\GRAPHICS{0.8}{0.8}{GoogleColab/makeverify.png}
	\caption{Datei \FILE{Darknet} ohne Erweiterungen nach erfolgreicher Ausführung des make-Befehls erstellt}\label{ColabDarknet}
\end{figure}

\subsection{Hilfsfunktionen definieren}

Bei diesen drei Funktionen handelt es sich um Hilfsfunktionen, die es uns ermöglichen, das Bild in Ihrem Colab Notebook anzuzeigen, nachdem wir unsere Erkennungen durchgeführt haben, sowie Bilder in unsere Cloud \ac{vm} hoch- und herunterzuladen. \ref{ColabTensorFlow:Complete}

\begin{code}[H]
	\lstinputlisting[language=Python, firstnumber=1]{../Code/JetsonNano/helperfunctions.py}
	\caption{Nützliche Funktionen zum Anzeigen von Bildern, Hochladen oder Herunterladen von Bildern aus der Cloud VM} \label{ColabTensorFlow:Complete}
\end{code}

\subsection{Verschieben der Datensätze in die Cloud VM}

Um eine benutzerdefinierte YOLOv4-Version zu erstellen, benötigen wir Folgendes:

\begin{itemize}
	\item Beschrifteter benutzerdefinierter Datensatz
	\item Benutzerdefinierte Datei \FILE{.cfg} 
	\item Datei \FILE{obj.data} und  Dateien \FILE{obj.names}
	\item Dateien \FILE{train.txt}  and \FILE{test.txt} 
\end{itemize}

Der Inhalt dieser Dateien wird im Folgenden näher erläutert.

Zunächst werden zwei ZIP-Ordner für Trainings- (DATEI{obj.zip}) und Validierungsdatensätze (DATEI{test.zip}) erstellt. 

Für das Training - Anzahl der Bilder

Für die Validierung - 
Diese sollten auf ein Laufwerk hochgeladen werden, das später in die Cloud-VM verschoben und in das Verzeichnis \glqq data\grqq{} entpackt werden kann.

\begin{figure}[H]
	\centering
	\GRAPHICS{0.8}{0.8}{GoogleColab/objandtest.png}
	\caption{Obj- und Test-Ordner mit Trainings- bzw. Validierungsdatensätzen}
\end{figure}

Dies kann mit den folgenden Befehlen geschehen: 

Um beide Datensätze in das Stammverzeichnis der Colab-VM zu kopieren:

\medskip

\SHELL{!cp /mydrive/yolov4/obj.zip ../}

\medskip

\SHELL{!cp /mydrive/yolov4/test.zip ../}

\medskip

So entpacken Sie die Datensätze und ihren Inhalt, so dass sie sich jetzt im Ordner \glqq /darknet/data/\grqq{} befinden

\medskip

\SHELL{!unzip ../obj.zip -d data/}

\medskip

\SHELL{!unzip ../test.zip -d data/}

\subsection{Dateien für die Schulung konfigurieren}

In diesem Schritt müssen Sie Ihre benutzerdefinierten Dateien \FILE{.cfg}, \FILE{obj.data}, \FILE{obj.names}, \FILE{train.txt} und \FILE{test.txt} richtig konfigurieren.

Es ist wichtig, alle diese Dateien mit äußerster Vorsicht zu konfigurieren, da Tippfehler oder kleine Fehler zu großen Problemen mit dem benutzerdefinierten Training führen können.

Cfg-Datei: Zunächst wird die Datei \FILE{yolov4.cfg} auf das Google-Laufwerk kopiert und geändert. 

Die folgenden Änderungen werden vorgenommen: 

\begin{itemize}
  \item batch = 64 und subdivisions = 16
  \item max{\_}batches = 6000, steps = 4800, 5400 als die Anzahl der Klassen= 2. Normalerweise wird dies durch max{\_}batches = (\# of classes) * 2000 bestimmt, aber der Wert sollte nicht kleiner als 6000 sein. Daher habe ich mich für 6000 entschieden. Die Schritte werden durch Berechnung von (80 Prozent von max{\_}batches), (90 Prozent von max{\_}batches) bestimmt
  \item Die Anzahl der Filter sollte wie folgt geändert werden: filters = (\# of classes + 5) * 3. Da die Anzahl der Klassen = 2 ist, habe ich diesen Wert auf 21 gesetzt. 
  \item Die Anzahl der Klassen sollte in jeder \ac{yolo} Ebene. In der Datei \FILE{.cfg} gibt es 3 Yolo-Ebenen. Auch die Anzahl der Filter sollte in jeder Faltungsschicht direkt über jeder Yolo-Schicht eingestellt werden.
  \item learning{\_}rate=0.001
  \item width=416 (des Bildes)
  \item height=416 
\end{itemize}

Dateien \FILE{obj.names} und \FILE{obj.data}:
Diese Dateien sollten mit dem folgenden Inhalt erstellt werden. 
Die Datei \FILE{obj.names} sollte die Klassen in der Reihenfolge enthalten, die in der Originaldatei \FILE{label.txt} angegeben ist.

\begin{figure}[H]
	\centering
	\GRAPHICS{0.6}{0.6}{GoogleColab/obj1.png}
	\caption{Inhalt der Datei \FILE{obj.names}}
\end{figure}

Es sollte ein Sicherungsordner auf unserem Laufwerk erstellt werden, um die Gewichte zu speichern. Dieser wird 
Die Datei \FILE{obj.data} enthält die Pfade zum Backup, die Anzahl der Klassen, die Pfade zu den Dateien \FILE{train.txt} und \FILE{test.txt} usw., 
Die Erstellung von DATEI{train.txt} und DATEI{test.txt} wird weiter unten erläutert. 

\begin{figure}[H]
	\centering
	\GRAPHICS{0.8}{0.8}{GoogleColab/objdata.png}
	\caption{Inhalt der Datei \FILE{obj.data}}
\end{figure}

Die Dateien \FILE{obj.data} und \FILE{obj.names} sollten von unserem Laufwerk auf die Cloud-VM hochgeladen werden. 

\medskip

\SHELL{!cp /mydrive/yolov4/obj.names ./data} 

\SHELL{	!cp /mydrive/yolov4/obj.data ./data}

\medskip


Generating \FILE{train.txt} and \FILE{test.txt} : 

Die letzten Konfigurationsdateien, die benötigt werden, bevor wir mit dem Training unseres benutzerdefinierten Detektors beginnen können, sind die Dateien \FILE{train.txt} und \FILE{test.txt}, die die relativen Pfade zu allen unseren Trainings- und Validierungsbildern enthalten.\cite{TheAIGuy:2021}

\begin{code}[H]
	{\tiny
	\lstinputlisting[language=Python, firstnumber=1]{../Code/JetsonNano/GenerateTrain.py}
    }
	\caption{Skript zur Erstellung der Datei \FILE{train.txt} mit den Pfaden zu den Trainingsbildern}\label{TensorFlow:Complete}
\end{code}

\begin{code}[H]
	\lstinputlisting[language=Python, firstnumber=1]{../Code/JetsonNano/GenerateTest.py}
	\caption{Skript zur Erstellung der Datei \FILE{test.txt} mit den Pfaden zu den Testbildern}\label{TensorFlow:Complete}
\end{code}

Diese Skripte können mit den Befehlen ausgeführt werden:

\medskip

\SHELL{!python generate\_train.py}

\SHELL{!python generate\_test.py}


\subsection{Herunterladen der vortrainierten Gewichte für die Faltungsschichten}

In diesem Schritt werden die Gewichte für die Faltungsschichten des YOLO V4-Netzwerks heruntergeladen. Durch die Verwendung dieser Gewichte wird unser benutzerdefinierter Objektdetektor viel genauer und muss nicht so lange trainieren.

Der Befehl zum Herunterladen dieser Gewichte \cite{Bochkovskiy:2020}

\medskip

\SHELL{!wget https://github.com/AlexeyAB/darknet/releases/download/darknet\_yolo\_v3\_optimal/yolov4.conv.137}

\subsection{Training des benutzerdefinierten Objektdetektors}

Um das Modell zu trainieren, sollten wir den folgenden Befehl ausführen.

\medskip

\SHELL{!./darknet detector train data/obj.data cfg/yolov4-obj.cfg yolov4.conv.137 -dont\_show -map}

\medskip

\begin{figure}[H]
	\centering
	\GRAPHICS{0.8}{0.8}{GoogleColab/training.png}
	\caption{Training des Modells in Google Colab}
\end{figure}

\subsection{Herausforderungen}

\begin{itemize}
  \item Laufzeitbeschränkungen : Das Google Colab stoppt nach 2000 Iterationen und wir müssen die Google Colab VM nach jeder Unterbrechung neu einrichten, um das Modell erneut zu trainieren.
  \item Fehler aufgetreten, dass die \ac{gpu} nicht weiter verwendet werden kann.
  \item Viel Zeit für 2000 Iterationen benötigt ca. 6 Stunden.
\end{itemize}

\begin{figure}[H]
	\centering
	\GRAPHICS{0.8}{0.8}{GoogleColab/failed.png}
	\caption{Training scheiterte nach 2000 Iterationen aufgrund der Laufzeitbeschränkungen von Google Colab}
\end{figure}

\begin{figure}[H]
	\centering
	\GRAPHICS{0.8}{0.8}{GoogleColab/fail.png}
	\caption{Kein Zugriff auf die GPU möglich} 
\end{figure}
